\documentclass[10pt,a4paper]{amsart}
\usepackage[margin=19mm]{geometry}
\usepackage{amsmath,amssymb,mathtools}
\usepackage[scr=rsfs,cal=euler]{mathalfa}
\usepackage{xfrac}
\usepackage[unicode,hidelinks,bookmarks=false]{hyperref}
\newcommand{\ie}{i.e.\ }
\newcommand{\cf}{cf.\ }
\newcommand{\resp}{respectively\ }
\newcommand{\Subsection}{\subsection}
\providecommand{\nobreakdash}{\nobreak\hbox{-}}
\setlength{\emergencystretch}{2em}
\multlinegap0pt
\usepackage{bm}
\usepackage{algorithm}\usepackage{algpseudocode}
\usepackage{xcolor}
\newtheorem{theorem}{Theorem}
\newtheorem{lemma}{Lemma}
\usepackage{cleveref}
\crefname{theorem}{Theorem}{Theorems}
\crefname{lemma}{Lemma}{Lemmas}
\newcommand{\R}{\mathbb{R}}
\newcommand{\diam}{\mathrm{diam}}
\def\eqref#1{equation~\ref{#1}}
\newcommand*{\norm}[1]{\left\|#1\right\|}
\title{Towards Scalable Persistence-Based Topological Optimization}
\author{Abderrahim Bendahi, Alexandre Duplessis, Arnaud Fickinger}
\date{Source: arXiv:2605.10996v1 (CC BY 4.0)}
\begin{document}
\maketitle

\noindent\emph{Source and licence.} Towards Scalable Persistence-Based Topological Optimization. Version arXiv:2605.10996v1 (CC BY 4.0), distributed by arXiv under the Creative Commons Attribution 4.0 International licence, http://creativecommons.org/licenses/by/4.0/, as recorded in the arXiv licence metadata for this exact version and shown on the abstract page https://arxiv.org/abs/2605.10996v1. Full licence notice: \texttt{licenses/arxiv-2605.10996-CC-BY-4.0.txt}.\par\smallskip

\noindent\textit{Editorial scope.} Counted unit: the article's theorem thm:nw\_lipschitz (source0.tex lines 879--887). The article's complete original proof of it is delivered with it (source0.tex lines 1357--1417, one \texttt{proof} environment of 61 lines, which the article places away from the statement). No mathematics is written, repaired or re-derived: the statement and the proof are the article's own lines, in the article's own notation, and no retained line is substituted or re-flowed.  No further source-local context is needed: every cross-reference of every retained line resolves inside the two delivered spans.  Nothing was omitted from any retained span, so the package carries no omission record.  No retained line contains a citation, so the wrapper carries no bibliography and no bibliography entry.  No retained line contains a figure environment, an image-inclusion command or drawing code, so no image is delivered and the package declares no asset dependency.  Editorial formatting only, in addition: an AMS \texttt{amsart} preamble that restates the article's own declarations and macros used by the retained lines, namely the environments \texttt{theorem}, \texttt{lemma} on separate counters, together with the standard packages those lines need.  The retained environments print in this document as Theorem 1, Lemma 1 in order of appearance, whereas the article prints the same items with the numbering of its own full document.  The complete native source of the article as distributed in the arXiv e-print for this exact version is bundled unchanged as \texttt{sources/200331/source0.tex}.
\begin{theorem}[Global Lipschitz constant of the NW field]
\label{thm:nw_lipschitz}
Assume $\norm{g_i}_2\le M$ for all $i \in I$. Then $\bar v$ is globally Lipschitz and
\begin{equation}
\label{eq:nw_lipschitz}
    \|D\bar v(x)\|_{\mathrm{op}} \le \frac{\diam(P)}{\sigma^2} M \qquad \forall x \in \R^d.
\end{equation}
In particular, $\norm{\bar v(x) - \bar v(y)}_2 \le L \norm{x-y}_2$ for all $x, y$ with $L := \frac{\diam(P) M}{\sigma^2}$.
\end{theorem}

\begin{proof}[Proof of \Cref{thm:nw_lipschitz}]

We start by stating and proving the following lemma which provides an expression of the gradient of the weights and bounds its norm .



\begin{lemma}[Gradient of weights]
\label{lem:grad_weights}
Let $\bar x(x):=\sum_{j \in I} w_j(x) x_j$ be the weighted barycenter of anchors at $x$.
Then for all $i \in I$ and $x \in \R^d$,
\begin{equation} \label{eq:grad_w}
    \nabla w_i(x) = \frac{w_i(x)}{\sigma^2} \left( x_i - \bar x(x) \right).
\end{equation}
Moreover, for all $x$ and $i$,
    \[ \norm{\nabla w_i(x)}_2 \le \frac{\diam(P)}{\sigma^2}w_i(x), \]
where $\diam(P) := \max_{p,q\in I}\|x_p - x_q\|_2$
\end{lemma} 

\begin{proof}[Proof of \Cref{lem:grad_weights}]
Using 
    \[\phi_i(x) = \exp\left(- \dfrac{\norm{x - x_i}^2}{2 \sigma^2} \right) \quad \text{ and } \quad w_i= \dfrac{\phi_i}{Z}, \]

we have, 
    \[ \nabla\phi_i(x) = - ( x - x_i ) \frac{\phi_i(x)}{\sigma^2} , \] 
and
    \[\nabla Z(x) = \sum_j \nabla \phi_j(x) = - \dfrac{1}{\sigma^2}  \sum_j ( x -x_j ) \phi_j(x). \]
By the quotient rule,
\begin{align*}
   \nabla w_i(x)& =\dfrac{\nabla\phi_i(x) Z(x)-\phi_i(x)\nabla Z(x)}{Z(x)^2} \\ &= -\dfrac{\phi_i(x)}{\sigma^2 Z(x)}
\Big[(x-x_i)-\sum_j w_j(x)(x-x_j)\Big]. 
\end{align*}

Since 
\begin{align*}
    \sum_j w_j(x)(x-x_j) &= x-\sum_j w_j(x)x_j
    \\ &=x-\bar x(x),
\end{align*} 
we obtain \eqref{eq:grad_w}.


Hence, 
    \[ \|\nabla w_i(x)\|_2 = \dfrac{w_i(x)}{\sigma^2}\|x_i-\bar x(x)\|_2. \]
Since $\bar x(x)$ is a convex combination of points in $P$, it lies in $\mathrm{Conv}(P)$, hence 
    \[ \|x_i-\bar x(x)\|_2 \le \diam(P). \]

\end{proof}



Differentiate $\bar v(x)=\sum_i w_i(x)g_i$:
    \[ D \bar v(x)=\sum_{i\in I} ( \nabla w_i(x) )\otimes g_i.\]
Using $\norm{a \otimes b}_{\mathrm{op}}=\|a\|_2 \|b\|_2$,
    \begin{align*}
       \norm{D \bar v(x)}_{\mathrm{op}} &\le \sum_i \norm{\nabla w_i(x)}_2 \norm{g_i}_2 \\ &\le M \sum_i \norm{\nabla w_i(x)}_2.  
    \end{align*} 
Apply Lemma~\ref{lem:grad_weights}:
    \begin{align*}
        \sum_i \norm{\nabla w_i(x)}_2 &\le \dfrac{\diam(P)}{\sigma^2}\sum_i w_i(x) \\ &= \dfrac{\diam(P)}{\sigma^2}, 
    \end{align*}
giving \eqref{eq:nw_lipschitz}. The Lipschitz inequality follows from the mean value theorem.
\end{proof}

\end{document}
